% Figure 3: Six-framework taxonomy diagram
% Two-column layout, 3 frameworks per column, color-coded by originating reviewer.

\usepackage{tikz}
\usetikzlibrary{shapes.rounded,positioning,fit}

% Color definitions (kept local to avoid clashes in acmart)
\colorlet{fw-ahagreen}{green!45!black}
\colorlet{fw-loadblue}{blue!55!gray}
\colorlet{fw-worldpurple}{violet!70!blue}
\colorlet{fw-socialorange}{orange!80!red}
\colorlet{fw-percteal}{teal!80!black}
\colorlet{fw-visrose}{red!55!white}

\begin{figure}[tb]
\centering

\begin{tikzpicture}[
  font=\small,
  fwbox/.style={
    draw=#1,
    fill=#1!8,
    rounded corners=4pt,
    line width=0.9pt,
    text width=5.3cm,
    inner sep=6pt,
    align=left
  },
  colsep/.style={
    draw=black!25,
    line width=0.5pt,
    dashed
  }
]

% ---------------------------------------------------------------
% LEFT COLUMN  (frameworks 1, 2, 3)
% ---------------------------------------------------------------

\node[fwbox=fw-ahagreen] (fw1) at (0,0) {%
  {\bfseries\color{fw-ahagreen}A-ha Economy}\\[1pt]
  {\itshape\small Foggy Brume}\\[2pt]
  Pacing of insight moments: the solver's \mbox{``Eureka''} cadence should be earned and deliberate.\\[3pt]
  \textcolor{black!50}{\scriptsize Transfers to: level design, game design}
};

\node[fwbox=fw-loadblue, below=0.45cm of fw1] (fw2) {%
  {\bfseries\color{fw-loadblue}Load-Bearing Test}\\[1pt]
  {\itshape\small Thomas Snyder}\\[2pt]
  Every element must earn its place; remove anything that could vanish without consequence.\\[3pt]
  \textcolor{black!50}{\scriptsize Transfers to: level design, narrative design}
};

\node[fwbox=fw-worldpurple, below=0.45cm of fw2] (fw3) {%
  {\bfseries\color{fw-worldpurple}World-Model Consistency}\\[1pt]
  {\itshape\small Emily Short}\\[2pt]
  Puzzle logic must hold within its fiction; any rule violation breaks the world, not just the mechanic.\\[3pt]
  \textcolor{black!50}{\scriptsize Transfers to: interactive fiction, tabletop RPG}
};

% ---------------------------------------------------------------
% RIGHT COLUMN  (frameworks 4, 5, 6)
% ---------------------------------------------------------------

\node[fwbox=fw-socialorange, right=0.65cm of fw1] (fw4) {%
  {\bfseries\color{fw-socialorange}Social Fabric}\\[1pt]
  {\itshape\small Jordan Weisman}\\[2pt]
  Does the puzzle require or reward collaboration? Interaction structure is a design variable.\\[3pt]
  \textcolor{black!50}{\scriptsize Transfers to: tabletop design, escape rooms}
};

\node[fwbox=fw-percteal, below=0.45cm of fw4] (fw5) {%
  {\bfseries\color{fw-percteal}Perceptual-Shift}\\[1pt]
  {\itshape\small Scott Kim}\\[2pt]
  The answer reframes what came before; the final reveal restructures the solver's memory of the puzzle.\\[3pt]
  \textcolor{black!50}{\scriptsize Transfers to: visual design, interaction design}
};

\node[fwbox=fw-visrose, below=0.45cm of fw5] (fw6) {%
  {\bfseries\color{fw-visrose}Visual Grammar}\\[1pt]
  {\itshape\small Dana Young}\\[2pt]
  Layout communicates structure before content is read; spatial arrangement is a primary semantic channel.\\[3pt]
  \textcolor{black!50}{\scriptsize Transfers to: information design, accessible design}
};

% ---------------------------------------------------------------
% Column divider — vertical rule centered between the two columns
% ---------------------------------------------------------------
\draw[colsep]
  ($ (fw1.north east)!0.5!(fw4.north west) $)
  --
  ($ (fw3.south east)!0.5!(fw6.south west) $);

% Column header labels
\node[font=\scriptsize\bfseries\color{black!45}, above=0.18cm of fw1, anchor=south]
  {Frameworks 1--3};
\node[font=\scriptsize\bfseries\color{black!45}, above=0.18cm of fw4, anchor=south]
  {Frameworks 4--6};

\end{tikzpicture}

\caption{The six evaluative frameworks surfaced by C6 full-profile reviews but absent in
C0--C5 reviews. Each framework is tied to a specific expert identity; the originating
reviewer introduced the framework's vocabulary across multiple reviews without prompting.
All six transfer to at least one creative domain beyond puzzle hunt design.}
\label{fig:taxonomy}
\end{figure}
